\section{Conclusion}
\label{sec:conclusion}

This paper has presented \software{}, a local-first, human-in-the-loop
platform that reframes AI-assisted academic writing as a document engineering
problem. By grounding AI assistance in a filesystem-backed project model with
permission-aware interaction, typed workflow pipelines, and multi-API citation
verification, the system addresses three practical problems---context
fragmentation, unaudited edits, and unverifiable artifacts---that general-purpose
chat interfaces leave unresolved.

The implementation comprises approximately 23,100 lines of application source
code across a monorepo with a React/Vite frontend and Fastify/Node.js backend,
validated by a 32-file Vitest suite containing 227 automated tests and a
Playwright end-to-end scenario.

Our experience building and using \software{} suggests three broader lessons
for the design of AI-assisted writing tools:

\begin{enumerate}
  \item \textbf{Project-level awareness of manuscript structure enables
  capabilities that text-level tools cannot provide.} When the system
  understands chapters, citations, and compilation dependencies, it can
  provide context-aware assistance, verify generated references, and validate
  edits through compilation---all without additional user effort.

  \item \textbf{Explicit permission boundaries increase user trust and
  encourage more exploratory AI use.} Separating AI interaction into Chat,
  Agent, and Tools modes creates clear mental models and reduces anxiety about
  unintended changes. Users who were reluctant to use AI in ``auto-edit''
  modes became comfortable with the gated approval workflow.

  \item \textbf{Local-first architectures provide practical advantages in
  tool compatibility, data ownership, and configuration flexibility.}
  Manuscripts remain ordinary git repositories, external tools work
  naturally, and different projects can use different LLM providers without
  application changes.
\end{enumerate}

Three directions merit immediate attention. First, controlled user studies
comparing document-engineering-aware AI assistance against chat-only or
Overleaf-only workflows would provide the quantitative evidence that our
current evaluation lacks. Second, scaling the filesystem-backed model to
institutional deployments with hundreds of concurrent users requires a metadata
indexing layer while preserving local-first benefits. Third, LLM-based compile
error repair---feeding compilation logs back to the model for automatic
correction---would close the loop between verification and remediation.

\software{} is available as open-source software under the MIT license.
The source code, documentation, and demonstration projects are accessible at
\repo{}. We invite contributions from the research software community and look
forward to collaborating on controlled user studies to further validate the
design claims in this paper.

% ===========================================================================
% DATA AVAILABILITY
% ===========================================================================
\section*{Data Availability}

The software repository includes a minimal redistributable demonstration paper
project under \texttt{examples/demo-paper/}. No private manuscript data or user
data is included. Performance data reported in this paper (citation
verification latency, compilation times) were collected from the production
APIs and development environment described in Section~\ref{sec:evaluation} and
are reproducible using the open-source codebase.

% ===========================================================================
% ACKNOWLEDGEMENTS
% ===========================================================================
\section*{Acknowledgements}

The authors thank the maintainers of the open-source tools and scholarly
infrastructure used by \software{}, including Node.js, React, Vite, Fastify,
CodeMirror, KaTeX, \TeX{} Live, CrossRef, Semantic Scholar, OpenAlex,
and the Model Context Protocol communities.

We also thank the early adopters who provided valuable feedback during the
development of \software{}, particularly those who tested the permission mode
system and pipeline workflows and contributed usability insights that shaped
the final design.

% ===========================================================================
% FUNDING
% ===========================================================================
\section*{Funding}

The author received no specific grant from any funding agency in the public,
commercial, or not-for-profit sectors to support this work.

% ===========================================================================
% COMPETING INTERESTS
% ===========================================================================
\section*{Competing Interests}

The author declares no competing interests related to this work.

% ===========================================================================
% CRediT AUTHOR CONTRIBUTION
% ===========================================================================
\section*{Author Contributions}

\textbf{ZuKang Xu:} Conceptualization, Software, Investigation, Writing ---
original draft, Writing --- review \& editing, Visualization.

% ===========================================================================
% AI ASSISTANCE DISCLOSURE (Wiley policy requirement)
% ===========================================================================
\section*{AI Assistance Disclosure}

During the preparation of this manuscript, the author used large language
model (LLM) tools (OpenAI GPT-4o and Anthropic Claude) for language
polishing, grammar correction, and rephrasing of selected passages. The
author has reviewed and edited all AI-generated content and takes full
responsibility for the final manuscript. These tools were not used for
generating research ideas, experimental designs, data analysis, or core
argumentation. No AI tool is listed as an author.
